\section{Line covers after random deletions}\label{sec:linecover}

We regard the affine plane $\F_q^2$ as a subset of its projective
completion $\PP^2(\F_q)$. Thus a pencil with centre at infinity is a
parallel class of affine lines. The purpose of this section is to show
that a set of at most $q$ lines covering almost all survivors of a
suitable random deletion process must contain almost a full pencil.
The deletions within a single trial may be dependent; independence is
required only between trials.

Classical covering results for finite planes are closely related by
duality to blocking sets. Jamison and Brouwer--Schrijver established the
affine blocking-set lower bound $2q-1$~\cite{Jamison1977,BrouwerSchrijver1978},
and Blokhuis, Brouwer, and Sz\H{o}nyi studied partial projective covers
and bounds depending on their uncovered points
\cite[Section 1.1, Proposition 1.5]{BBS2010}. Here the covered set instead consists of the
positive-density survivors of structured random deletions. We prove the
pencil statement needed in this setting directly.

There are two different concentration tasks. The total number of
surviving points and the number on a fixed line have bounded changes
under one trial; for these we use the elementary estimates
in~\cite[Lemmas~A.2--A.3]{OpenAI2026Balanced}. The more delicate task
is to retain many intersections of a short list of lines, simultaneously
for every such list. A deletion trial rarely hits this small set, and
we must use that rarity in the variance estimate. We prove the needed
Bernstein-type martingale bound here, following the method of
Freedman~\cite[Theorem (1.6) and Proposition (2.1)]{Freedman1975}.

\begin{lemma}[Martingale concentration]\label{lem:martingale}
Let $Z_i$, $1\leq i\leq n$, be martingale differences for a filtration
$(\mathcal A_i)_{i=0}^n$. Suppose that $M,V>0$ and deterministic numbers
$v_i\geq0$ satisfy, almost surely,
\[
 |Z_i|\leq M,\qquad
 \E(Z_i^2\mid\mathcal A_{i-1})\leq v_i,\qquad
 \sum_{i=1}^n v_i\leq V.
\]
Then, for every $u>0$,
\begin{equation}\label{eq:lc-martingale}
 \Pr\left(\left|\sum_{i=1}^n Z_i\right|\geq u\right)
 \leq 2\exp\left[-\frac14
   \min\left\{\frac{u^2}{V},\frac{u}{M}\right\}\right].
\end{equation}
\end{lemma}

\begin{proof}
For $|z|\leq1$ we have $e^z\leq1+z+z^2$. Consequently, if
$|\theta|\leq M^{-1}$, the conditional mean-zero property gives
\[
 \E(e^{\theta Z_i}\mid\mathcal A_{i-1})
 \leq 1+\theta^2v_i\leq e^{\theta^2v_i}.
\]
Iteration and Markov's inequality bound the upper tail by
$\exp(\theta^2V-\theta u)$ for $0<\theta\leq M^{-1}$.
Take $\theta=\min\{u/(2V),1/M\}$. If $u\leq2V/M$, the exponent is
$-u^2/(4V)$; otherwise it is at most $-u/(2M)$.
Applying the same argument to $-Z_i$ proves
\eqref{eq:lc-martingale}.

\end{proof}

\begin{lemma}[Line covers after random deletions]\label{lem:linecover}
Let $q\to\infty$ through prime powers for which every proper subfield
of $\F_q$ has order at most $q^{1/3}$. Fix a positive integer $D$ and
constants $\delta>0$ and $0<\rho<1$. For each $q$, let $\mathcal F$ be a
family of at most $Dq^D$ subsets of $\F_q^2$, called shapes, such that
\begin{enumerate}
 \item every $F\in\mathcal F$ has at most $Dq$ points;
 \item for every $F\in\mathcal F$ and affine line $\ell$, either
       $F=\ell$ or $|F\cap\ell|\leq D$.
\end{enumerate}
There are at most $Dq$ independent trials, each producing a random list
of at most $D$ shapes from $\mathcal F$. Delete every point on any listed
shape, and denote the surviving set by $X$. Assume that
\begin{enumerate}
 \item for each point $x$ and each trial, the probability that this
       trial deletes $x$ is at most $D/q$;
 \item $\Pr(x\in X)=\delta_q\geq\delta$ for every $x\in\F_q^2$;
 \item for each affine line $\ell$, the sum over trials of the
       probabilities that the trial lists the exact shape $\ell$ is
       at most $D/q$.
\end{enumerate}
Write $X^\ell$ for the surviving set when all exact occurrences of the
shape $\ell$ are ignored, with every other listed shape still deleted.
With probability tending to one, both of the following statements
hold:
\begin{enumerate}
 \item $|X^\ell\cap\ell|\geq\delta q/2$ for every affine line $\ell$;
 \item every set of at most $q$ affine lines covering all but at most
       $2q$ points of $X$ contains at least $(1-\rho)q$ lines through a
       common projective point.
\end{enumerate}
\end{lemma}

\begin{proof}
We first establish a simultaneous lower bound on surviving
cross-intersections for every admissible pair of short line lists,
ignoring any deletion shape equal to one of those lines. After fixing
such a realization, we show that any cover lacking an almost full
pencil contains an admissible pair with too few surviving
cross-intersections. This second step samples only within the fixed
cover, with separate arguments according to its maximum concurrency.

Fix the absolute constant $\kappa\in(0,1/10)$ of
Theorem~\ref{thm:incidence}, and put $\epsilon=\kappa/10$.
All constants below may depend on the fixed parameters, but not on
$q$, the particular deletion distributions, or a subsequently chosen
cover.

\paragraph{Uniform point and line counts.}
Expose the independent trials one at a time. Changing one trial
changes $|X|$ by at most $2D^2q$, because the old and new lists contain
at most $2D$ shapes in total and each has at most $Dq$ points.
The exposure observation in~\cite[Lemma~A.3]{OpenAI2026Balanced}
therefore bounds the absolute martingale increment by $L=2D^2q$.
There are $T\le Dq$ trials, and $\E|X|=\delta_q q^2$. The fixed-time
concentration estimate in~\cite[Lemma~A.2]{OpenAI2026Balanced}, with
$u=q^{3/2+\epsilon}$, gives
\[
 \Pr\bigl(\bigl||X|-\delta_q q^2\bigr|\ge q^{3/2+\epsilon}\bigr)
 \le2\exp\left(-\frac{q^{2\epsilon}}{32D^5}\right).
\]
If there are no trials, the count is deterministic and this bound
holds as well. Equivalently, we may add a dummy trial of zero change.

For the line count $|X^\ell\cap\ell|$, the rule ignoring exact copies
of $\ell$ is applied separately to each trial. Changing one trial
then changes the count by at most $2D^2$: each of its at most $2D$
old or new shapes meets $\ell$ in at most $D$ points. With
$L=2D^2$ and $u=q^{1/2+\epsilon}$, the same two elementary estimates
and $T\le Dq$ give
\[
 \Pr\bigl(\bigl||X^\ell\cap\ell|-\E|X^\ell\cap\ell|\bigr|
              \ge q^{1/2+\epsilon}\bigr)
 \le2\exp\left(-\frac{q^{2\epsilon}}{32D^5}\right).
\]

Moreover, a point restored by ignoring $\ell$ can be restored only if
some trial lists $\ell$ exactly. The probability of this event is at
most $D/q$, by the third deletion hypothesis. Therefore
\[
 \delta_q q\leq \E|X^\ell\cap\ell|\leq\delta_q q+D.
\]
There are $q(q+1)$ affine lines. The superpolynomial tails and a union
bound consequently show that, with probability tending to one,
\begin{equation}\label{eq:lc-counts}
 \begin{split}
 |X|&=\delta_q q^2+O_D(q^{3/2+\epsilon}),\\
 |X^\ell\cap\ell|&=\delta_q q+O_D(q^{1/2+\epsilon})
       \qquad\text{for every affine line }\ell.
 \end{split}
\end{equation}
The lower line estimate proves the first conclusion of the lemma.
For the pencil conclusion we will also need to bound overlaps in an
alleged cover. The same estimates give the uniform upper bound
\begin{equation}\label{eq:lc-line-upper}
 |X\cap\ell|\leq\delta_q q+O_D(q^{1/2+\epsilon}),
\end{equation}
because $X\subseteq X^\ell$.

\paragraph{A simultaneous event for small line samples.}
The counts just proved control a single line but do not yet restrict
how a cover's lines meet. We now obtain a stronger test on short lists
of lines: provided their intersections are distinct and not concentrated
on a deletion shape, many of those intersections survive. The list
length will be logarithmic in $q$, so that the variance-sensitive tail
can pay for a union bound over all lists.

Set
\[
 h=10(D+1),\qquad H_0=2Dh,\qquad
 t=\lceil A\log q\rceil,
\]
where the fixed constant $A$ will be chosen below. Consider any ordered
tuple of $2t$ distinct affine lines, separated into two lists of length
$t$. Call the tuple admissible if its cross-list intersections form a
set $P_0$ of $t^2$ distinct finite points, and
\begin{equation}\label{eq:lc-occupancy}
 |F\cap P_0|\leq H_0
 \quad\text{for every }F\in\mathcal F
       \text{ which is not one of the tuple lines}.
\end{equation}
For this tuple, ignore deletions by every shape equal to a tuple line,
and let $Y_0$ count the survivors in $P_0$ after the remaining deletions.
We claim that, with probability tending to one,
\begin{equation}\label{eq:lc-tuple-event}
 Y_0\geq\frac{\delta t^2}{2}
 \qquad\text{for every admissible tuple}.
\end{equation}

To prove the claim, fix a tuple in advance. Omitting shapes is a
deterministic operation on each trial, so the trials remain
independent. Omission can only restore points, and hence
$\E Y_0\geq\delta t^2$. By \eqref{eq:lc-occupancy}, each trial's
contribution to the deletion set within $P_0$ has size at most
$C_0=DH_0$. The probability that this contribution is nonempty is at
most $Dt^2/q$, by a union bound over $P_0$ and the per-point deletion
hypothesis.

For completeness, the rarity of a nonempty contribution also controls
the conditional variance of the exposure martingale. Fix the past,
and let $g_i(w)$ be the conditional expected final value of $Y_0$ when
the current trial has outcome $w$. Let $g_i(\varnothing)$ denote the
hypothetical value obtained by making this trial delete nothing on
$P_0$. Independence leaves the law of the future trials unchanged, so
\[
 |g_i(w)-g_i(\varnothing)|\leq C_0,
\]
with equality to zero whenever the trial's contribution is empty.
The current trial also retains its original law conditional on the
past. Thus the martingale increment has absolute value at most $2C_0$, and
its conditional variance is at most
\[
 \E\bigl((g_i(w)-g_i(\varnothing))^2\mid\text{past}\bigr)
 \leq C_0^2\frac{Dt^2}{q}.
\]
The sum of conditional variances is at most $D^2C_0^2t^2$, since there
are at most $Dq$ trials. Applying Lemma~\ref{lem:martingale} at deviation
$\delta t^2/2$ yields
\begin{equation}\label{eq:lc-tuple-tail}
 \Pr\left(Y_0<\frac{\delta t^2}{2}\right)
 \leq2e^{-c_0t^2},\qquad
 c_0=\min\left\{\frac{\delta^2}{16D^2C_0^2},
                       \frac{\delta}{16C_0}\right\}>0.
\end{equation}
In particular, $c_0$ is independent of $A$.

There are at most $(q(q+1))^{2t}$ ordered tuples. The logarithm of the
union bound in \eqref{eq:lc-tuple-tail} is
\[
 -c_0t^2+4t\log q+O(t/q)+O(1)
 =\bigl(-c_0A^2+4A+o(1)\bigr)(\log q)^2.
\]
Choose, for example, any fixed $A>5/c_0$. This proves
\eqref{eq:lc-tuple-event} simultaneously for all admissible tuples.

\paragraph{Consequences of an alleged cover.}
The random part of the argument is complete. We now keep the deletion
outcome fixed and show deterministically that the simultaneous events
exclude any cover without an almost full pencil. Random sampling below
is only a way to find a short witness inside a fixed cover.

Fix a realization satisfying both \eqref{eq:lc-counts} and
\eqref{eq:lc-tuple-event}. Suppose that the second conclusion of the
lemma fails. From an offending cover remove every line which actually
occurred as a listed deletion shape. Such a line contains no point of
$X$, so the remaining line set $L$ still covers all but at most $2q$
points of $X$. Put $m=|L|$, and for any projective point $p$ write
$n_p=|\{\ell\in L:p\in\ell\}|$. Define
\[
 E_L=\sum_{x\in X}(n_x-1)_+,
 \qquad (a)_+=\max\{a,0\}.
\]
The count of covered survivors and \eqref{eq:lc-line-upper} give
\[
 m\bigl(\delta_q q+O_D(q^{1/2+\epsilon})\bigr)
 \geq\delta_q q^2-O_D(q^{3/2+\epsilon})-2q.
\]
Since $\delta_q\geq\delta$ and $m\leq q$, this implies
\begin{equation}\label{eq:lc-cover-counts}
 m=q-O_{D,\delta}(q^{1/2+\epsilon}),\qquad
 E_L=O_D(q^{3/2+\epsilon}).
\end{equation}
For the second estimate, subtract the number of covered survivors
from $\sum_{\ell\in L}|X\cap\ell|$ and use $m\leq q$.
Finally, put
\[
 R=\max_{p\in\PP^2(\F_q)}n_p.
\]
The failure of the pencil conclusion implies $R\leq(1-\rho)q$.
We will find an admissible tuple of lines from $L$ with fewer than
$\delta t^2/2$ survivors in its cross-intersection set.

\paragraph{Low concurrency: $R\leq q^{1/2+\kappa/4}$.}
Let $T=q^{1/2-\kappa/3}$. The number of ordered pairs of distinct lines
of $L$ meeting at a point $x\in X$ is $\sum_{x\in X}n_x(n_x-1)$.
Points with $n_x<T$ contribute at most
\begin{equation}\label{eq:lc-low-light}
 T E_L=O(q^{2-7\kappa/30}).
\end{equation}
The set $P=\{x\in X:n_x\geq T\}$ has size
\[
 |P|\leq\frac{E_L}{T-1}
      =O(q^{1+13\kappa/30})\leq q^{1+\kappa}
\]
for sufficiently large $q$. Theorem~\ref{thm:incidence}, applied to
$P$ and $L$, therefore gives $\sum_{x\in P}n_x\leq q^{3/2-\kappa}$.
The contribution of these points to the ordered-pair count is at most
\begin{equation}\label{eq:lc-low-heavy}
 R\sum_{x\in P}n_x\leq q^{2-3\kappa/4}.
\end{equation}
Together, \eqref{eq:lc-low-light} and \eqref{eq:lc-low-heavy} show that
the fraction of ordered distinct line pairs meeting in $X$ is
\begin{equation}\label{eq:lc-low-pair-fraction}
 O\bigl(q^{-7\kappa/30}+q^{-3\kappa/4}\bigr)=o(1).
\end{equation}

Sample an ordered tuple of $2t$ distinct lines uniformly from $L$.
The expected number of its cross pairs meeting in $X$ is $t^2$ times
the fraction in \eqref{eq:lc-low-pair-fraction}. Markov's inequality
shows that this number is less than $\delta t^2/2$ with probability
tending to one.

We next check the two admissibility conditions. A given sampled pair
is parallel with probability at most $R/(m-1)$. Once two distinct
lines have been sampled, at most $R$ population lines pass through
their projective intersection, so the probability of concurrency for
any specified three sampled positions is at most $R/(m-2)$. The
probability of any parallel pair or concurrent triple in the tuple
is consequently
\begin{equation}\label{eq:lc-low-degeneracy}
 O\left(\frac{(t^2+t^3)R}{m}\right)
 =O\bigl((\log q)^3q^{-1/2+\kappa/4}\bigr)=o(1).
\end{equation}
If neither event occurs, every cross intersection is finite and the
cross intersections are distinct: equality of two different cross
intersections would make at least three of the tuple lines concurrent.

To verify \eqref{eq:lc-occupancy}, fix $F\in\mathcal F$. Consider the
event that $F$ is not a tuple line but contains more than $H_0$ distinct
finite cross intersections. Each tuple line then contains at most
$D$ points of $F$. Greedily choose a cross intersection on $F$ and its
two representing positions, and discard all cross intersections on
either of those lines. Each step discards at most $2D$ points. Since
$H_0=2Dh$, this produces $h$ cross pairs whose $2h$ positions are
distinct.

Among the population $L$, there are at most $mDR$ ordered distinct
line pairs, neither line equal to $F$, whose intersection lies in
$F$. Indeed, the first line meets $F$ in at most $D$ points, and at
each such point there are at most $R$ choices of the second line.
There are at most $t^{2h}$ ordered choices of the $h$ cross pairs of
positions. For each choice, the sampled lines at those $2h$ distinct
positions are uniform among $(m)_{2h}$ assignments, where
$(z)_a=z(z-1)\cdots(z-a+1)$. The probability of the event for this
fixed shape is therefore at most
\begin{equation}\label{eq:lc-low-shape}
 \frac{t^{2h}(mDR)^h}{(m)_{2h}}
 =O_{D,h}\left(\left(\frac{t^2R}{m}\right)^h\right).
\end{equation}
The numerator may count assignments with repeated population lines;
this only enlarges the upper bound. Summing
\eqref{eq:lc-low-shape} over all shapes gives
\[
 O\bigl((\log q)^{2h}q^{D-h(1/2-\kappa/4)}\bigr)=o(1),
\]
because $\kappa<1/10$ and $h=10(D+1)$ imply
\[
 D-h(1/2-\kappa/4)
 <D-\frac{19}{4}(D+1)<0.
\]
Thus, with probability tending to one, the sampled tuple is admissible
and has fewer than $\delta t^2/2$ cross pairs meeting in $X$.

\paragraph{High concurrency: $R>q^{1/2+\kappa/4}$.}
Choose a projective point $p$ with $n_p=R$, and take all $R$ lines of
$L$ through $p$ as the pencil. Its complement has size
\[
 N=m-R\geq\rho q/2
\]
for sufficiently large $q$, by \eqref{eq:lc-cover-counts}; no
complement line passes through $p$. Sample the first list of $t$ lines
uniformly without replacement from the pencil, and independently
sample the second list in the same way from the complement.

Any cross intersection lies off $p$. At a cross intersection
$x\in X$, exactly one pencil line passes through $x$, and the number
of cross pairs meeting there is the number of complement lines
through $x$, namely $n_x-1$. Thus at most $E_L$ of all $RN$ population
cross pairs meet in $X$. Their fraction is at most
\begin{equation}\label{eq:lc-high-pair-fraction}
 \frac{E_L}{RN}
 =O\bigl(q^{\epsilon-\kappa/4}\bigr)
 =O(q^{-3\kappa/20})=o(1).
\end{equation}
As before, expectation and Markov's inequality show that fewer than
$\delta t^2/2$ sampled cross pairs meet in $X$, with probability
tending to one.

For a fixed complement line, at most one pencil line meets it at
infinity when $p$ is finite. When $p$ is at infinity, the pencil lines
have direction $p$ and every complement line has a different
direction, so all cross intersections are finite. The probability of
any infinite cross intersection is therefore at most $t^2/R$.
Two equal cross intersections must use the same pencil line, since
distinct pencil lines meet only at $p$, which lies on no complement
line. Given two complement lines, their projective intersection
determines at most one affine pencil line, namely its join with $p$.
If this join is the line at infinity, there is no affine candidate.
Conditioning on the second list and taking a union bound over its
pairs shows that the probability of a collision is $O(t^3/R)$.
The total probability of a failure of finiteness or distinctness is
therefore
\begin{equation}\label{eq:lc-high-degeneracy}
 O\left(\frac{t^2+t^3}{R}\right)
 =O\bigl((\log q)^3q^{-1/2-\kappa/4}\bigr)=o(1).
\end{equation}

Fix $F\in\mathcal F$, and consider the event that $F$ is not a tuple
line but contains more than $H_0$ distinct finite cross intersections.
The same greedy argument gives $h$ cross pairs with disjoint
positions. There are at most
$DN$ population cross pairs, with neither line equal to $F$, whose
intersection lies in $F$: a complement line distinct from $F$ meets
$F$ in at most $D$ points, all off $p$, and each such point determines
at most one pencil line. The two lists are sampled independently, so
the probability of this occupancy violation for a fixed shape is at
most
\begin{equation}\label{eq:lc-high-shape}
 \frac{t^{2h}(DN)^h}{(R)_h(N)_h}
 =O_{D,h}\left(\left(\frac{t^2}{R}\right)^h\right).
\end{equation}
After summing over $\mathcal F$, the bound becomes
\[
 O\bigl((\log q)^{2h}q^{D-h(1/2+\kappa/4)}\bigr)=o(1),
\]
since
\[
 D-h(1/2+\kappa/4)<D-5(D+1)<0.
\]
The sampled tuple is consequently admissible and has fewer than
$\delta t^2/2$ cross pairs meeting in $X$, with probability tending
to one in this regime as well.

\paragraph{The contradiction and the order of choices.}
All sampling estimates above are uniform over line sets satisfying
\eqref{eq:lc-cover-counts} and the appropriate bound on $R$. In either
regime, for every sufficiently large $q$ there is therefore an
admissible tuple from the fixed line set $L$ with
\[
 |X\cap P_0|<\delta t^2/2.
\]
No line of $L$ actually occurred as a listed deletion shape. Hence
ignoring exact occurrences of its tuple lines changes none of the
deletions in the fixed realization, and the tuple's count $Y_0$ is
exactly $|X\cap P_0|$. This contradicts
\eqref{eq:lc-tuple-event}.

In particular, the cover and the tuple are allowed to depend on the
realized deletions: the event \eqref{eq:lc-tuple-event} was established
simultaneously for every admissible tuple before fixing the
realization. Sampling from $L$ is used only to prove the existence of
a tuple inside that fixed realization. No concentration estimate is
conditioned on the choice of the cover. The contradiction proves the
second conclusion, and finishes the proof.
\end{proof}
